\section{Effective complete-source moment tails}\label{sec:tails}

The moments used in the event chart retain the entire theta source and
the entire integration half-line. This appendix gives explicit upper
bounds for replacing them by finitely many modes on a finite interval.
For $S\ge0$, $Y\ge0$, and an integer $j\ge0$, recall the notation
\begin{equation}\label{eq:tails-moment}
 M_{j,Y}(S)=\int_0^\infty u^j e^{Su^2+Yu}\Phi(u)\dd u.
\end{equation}
Write $\lambda_n=\pi n^2$ and split the source in
\eqref{eq:Phi} into its complete individual summands,
\begin{equation}\label{eq:tails-summand}
 \Phi_n(u)=\lambda_n e^{5u}(2\lambda_ne^{4u}-3)
                         e^{-\lambda_ne^{4u}}.
\end{equation}
For $u\ge0$, the factor $2\lambda_ne^{4u}-3$ is at least
$2\pi-3>0$. Thus every $\Phi_n$ is positive, and
\begin{equation}\label{eq:tails-leading}
 0<\Phi_n(u)\le 2\lambda_n^2e^{9u-\lambda_ne^{4u}}.
\end{equation}
The smaller term in \eqref{eq:tails-summand} is discarded only in
this upper estimate; it remains part of the finite source below.

Choose an integer $N\ge0$ and a real $U\ge1$. Put $k=N+1$ and define
the exact finite head
\begin{equation}\label{eq:tails-head}
 M^{N,U}_{j,Y}(S)=\sum_{n=1}^{N}
       \int_0^U u^j e^{Su^2+Yu}\Phi_n(u)\dd u.
\end{equation}
An empty sum is understood as zero. Positivity and Tonelli's theorem,
initially allowing infinite values, give the disjoint decomposition
\begin{align}
 M_{j,Y}(S)=M^{N,U}_{j,Y}(S)
 &+\sum_{n=k}^\infty\int_0^\infty
                  u^j e^{Su^2+Yu}\Phi_n(u)\dd u\notag\\
 &+\sum_{n=1}^{N}\int_U^\infty
                  u^j e^{Su^2+Yu}\Phi_n(u)\dd u.
 \label{eq:tails-decomposition}
\end{align}
The first omitted piece comprises all omitted modes on the whole
half-line. The second comprises only retained modes on their exterior
rays. Their index--integration regions do not overlap.

\begin{proposition}[All omitted modes]\label{prop:tails-modes}
Set
\begin{equation}\label{eq:tails-mode-parameters}
 \rho_k=\left(\frac{k+1}{k}\right)^4e^{-\pi(2k+1)},
 \qquad b_k=4\lambda_k-9-Y.
\end{equation}
If $b_k>0$ and $S\le8\lambda_k$, then
\begin{equation}\label{eq:tails-modes}
 \sum_{n=k}^\infty\int_0^\infty
       u^j e^{Su^2+Yu}\Phi_n(u)\dd u
 \le L_{j,Y}(S;N):=
 \frac{2\lambda_k^2e^{-\lambda_k}j!}
                  {(1-\rho_k)b_k^{j+1}}.
\end{equation}
In particular $\rho_k<1$ for every $k\ge1$.
\end{proposition}

\begin{proof}
For $n\ge k$ and $u\ge0$, the ratio of consecutive leading terms in
\eqref{eq:tails-leading} is
\[
 \frac{\lambda_{n+1}^2e^{-\lambda_{n+1}e^{4u}}}
      {\lambda_n^2e^{-\lambda_ne^{4u}}}
 =\left(\frac{n+1}{n}\right)^4
                 e^{-\pi(2n+1)e^{4u}}
 \le\rho_k.
\]
Both the rational factor and the exponential factor are bounded above
by their values at $n=k$, $u=0$. Also
$\rho_k\le16e^{-3\pi}<1$. Geometric summation therefore gives
\begin{equation}\label{eq:tails-geometric}
 \sum_{n=k}^\infty\Phi_n(u)
 \le\frac{2\lambda_k^2}{1-\rho_k}
                 e^{9u-\lambda_ke^{4u}}.
\end{equation}
For $u\ge0$, Taylor's inequality for the positive exponential series is
\begin{equation}\label{eq:tails-exp-inequality}
 e^{4u}\ge1+4u+8u^2.
\end{equation}
Consequently
\begin{align*}
 Su^2+(9+Y)u-\lambda_ke^{4u}
 &\le-\lambda_k-(4\lambda_k-9-Y)u
                         +(S-8\lambda_k)u^2\\
 &\le-\lambda_k-b_ku.
\end{align*}
Multiply \eqref{eq:tails-geometric} by $u^je^{Su^2+Yu}$ and
integrate. The elementary integral
\[
 \int_0^\infty u^je^{-bu}\dd u=\frac{j!}{b^{j+1}}
 \qquad(b>0)
\]
follows by scaling and $j$ integrations by parts. It gives exactly
\eqref{eq:tails-modes}.
\end{proof}

\begin{proposition}[Retained exterior rays]\label{prop:tails-rays}
For each retained index $1\le n\le N$, suppose that
\begin{equation}\label{eq:tails-ray-parameters}
 b_n(U)=4\lambda_ne^{4U}-9-2SU-Y>0,
 \qquad S\le8\lambda_ne^{4U}.
\end{equation}
Then
\begin{align}
 \int_U^\infty u^je^{Su^2+Yu}\Phi_n(u)\dd u
 &\le R_{n,j,Y}(S;U),\label{eq:tails-ray-bound}\\
 R_{n,j,Y}(S;U)
 &:={2\lambda_n^2}
 e^{SU^2+(9+Y)U-\lambda_ne^{4U}}
 \sum_{h=0}^{j}\binom{j}{h}
       \frac{U^{j-h}h!}{b_n(U)^{h+1}}.
 \label{eq:tails-ray}
\end{align}
\end{proposition}

\begin{proof}
Set $u=U+v$, where $v\ge0$. Applying
\eqref{eq:tails-exp-inequality} to $v$ gives
\begin{align*}
 S(U+v)^2+(9+Y)(U+v)-\lambda_ne^{4U}e^{4v}
 \le{}&SU^2+(9+Y)U-\lambda_ne^{4U}\\
 &-b_n(U)v+(S-8\lambda_ne^{4U})v^2.
\end{align*}
The quadratic coefficient is nonpositive by hypothesis. Thus
\eqref{eq:tails-leading} bounds the exterior integral by
\[
 2\lambda_n^2e^{SU^2+(9+Y)U-\lambda_ne^{4U}}
       \int_0^\infty(U+v)^je^{-b_n(U)v}\dd v.
\]
The binomial theorem and the same elementary exponential integrals give
\[
 \int_0^\infty(U+v)^je^{-bv}\dd v
 =\sum_{h=0}^{j}\binom{j}{h}
                  \frac{U^{j-h}h!}{b^{h+1}},
\]
which proves \eqref{eq:tails-ray-bound}--\eqref{eq:tails-ray}.
\end{proof}

\begin{corollary}[Complete finite enclosure]\label{cor:tails-complete}
Under the hypotheses of Propositions~\ref{prop:tails-modes} and
\ref{prop:tails-rays},
\begin{equation}\label{eq:tails-complete}
 M^{N,U}_{j,Y}(S)\le M_{j,Y}(S)
 \le M^{N,U}_{j,Y}(S)+L_{j,Y}(S;N)
                       +\sum_{n=1}^{N}R_{n,j,Y}(S;U).
\end{equation}
For each fixed $S,Y,j$, admissible $N,U$ exist and the sum of the
displayed tail bounds can be made arbitrarily small.
\end{corollary}

\begin{proof}
Insert the two upper bounds into the positive disjoint identity
\eqref{eq:tails-decomposition}. To verify the final assertion, first
increase $N$. Then $b_k=4\pi k^2-9-Y$ is positive,
$8\pi k^2\ge S$, $\rho_k\to0$, and the right-hand side of
\eqref{eq:tails-modes} tends to zero because of $e^{-\pi k^2}$.
With such an $N$ fixed, increase $U$. For the finitely many retained
indices, the term $e^{4U}$ dominates both $U$ and the constant $S$;
hence \eqref{eq:tails-ray-parameters} eventually holds simultaneously.
In \eqref{eq:tails-ray}, the negative exponent
$-\lambda_ne^{4U}$ dominates $SU^2+(9+Y)U$, while the finite sum has
only polynomial powers of $U$ and inverse positive powers of $b_n(U)$.
Thus every $R_{n,j,Y}(S;U)$ tends to zero. This proves the assertion
and also proves finiteness of \eqref{eq:tails-moment} directly.
\end{proof}

These bounds estimate the complete positive moments. They do not impose
separate reflection identities on a source head or tail, and they make
no replacement of the canonical coefficient sequence. In particular, the
finite head in \eqref{eq:tails-head} retains both terms of every
summand in \eqref{eq:tails-summand}.
